% Part I, program calculation: ~13 min. Recursion schemes, their laws, a
% calculation by hand, homomorphisms, and our ICFP 2016 work
\partframe{Part I}{Program Calculation}

\begin{frame}[vertical-alignment = top]
  \frametitle{Programs as Algebra}
  \vspace*{4.5mm}
  Specify \emph{clearly}; \textbf{calculate} an efficient program;
  every step is an \emph{equation}.
  \par\medskip{\centering
    \begin{tikzpicture}[alt={A timeline of the algebra of programming.
        1978: Backus asks for an algebra of programs in his Turing lecture.
        1986 and 1987: Meertens and Bird, laws for lists. 1990 and 1991:
        Malcolm, and Meijer, Fokkinga and Paterson, laws for every data
        type. 1997: Bird and de Moor, the book Algebra of Programming. 2015:
        Hinze, Wu and Gibbons, one scheme behind them all. Forthcoming:
        Gibbons, Patterns in Functional Programming.},
      ev/.style = {box, font = \footnotesize, text width = 30mm,
        align = center, execute at begin node = {\hyphenpenalty = 10000}},
      yr/.style = {font = \footnotesize\bfseries, inner sep = 1pt}]
      \draw[very thick, gray, ->] (-0.6, 0) -- (14.6, 0);
      % \s = 1: event above the line, year below; \s = -1: the other way
      % each box: who and what, then a short citation (\tlcite)
      \foreach \x/\y/\s/\t [count = \i] in {
          0.65/1978/1/{\textbf{Backus}: ``an algebra of programs''
            \tlcite{CACM 1978}},
          3.15/1986--87/-1/{\textbf{Meertens}, \textbf{Bird}: laws for
            lists \tlcite{CWI'86; NATO ASI'87}},
          5.65/1990--91/1/{\textbf{Malcolm}; \textbf{Meijer} et al.: laws
            for every data type \tlcite{SCP 1990; FPCA 1991}},
          8.15/1997/-1/{\textbf{Bird}, \textbf{de Moor}: \emph{Algebra of
            Programming} \tlcite{Prentice Hall 1997}},
          10.65/2015/1/{\textbf{Hinze}, \textbf{Wu}, \textbf{Gibbons}: one
            scheme behind them all \tlcite{POPL 2015}},
          13.15/2027/-1/{\textbf{Gibbons}: \emph{Patterns in Functional
            Programming} \tlcite{CUP, forthcoming}}} {
        \ifnum\s=1 \def\ea{south}\def\ya{north}%
        \else \def\ea{north}\def\ya{south}\fi
        \fill[accent] (\x, 0) circle (1.2mm);
        \node[yr, anchor = \ya] at (\x, -0.25 * \s) {\y};
        \node[ev, anchor = \ea] (e\i) at (\x, 0.5 * \s) {\t};
        \draw[accent] (\x, 0.12 * \s) -- (e\i);
      }
    \end{tikzpicture}\par}
  \excerpt<2>{``Conventional programming languages are growing ever more
    enormous, but not stronger. [\ldots] Associated with the functional style
    of programming is an \textbf{algebra of programs} whose variables range
    over programs and whose operations are combining forms. This algebra can
    be used to \textbf{transform programs} and to solve equations whose
    `unknowns' are programs in much the same way one transforms equations in
    high school algebra. These transformations are given by algebraic laws
    [\ldots]''}{J. Backus, Can Programming Be Liberated from the von Neumann
    Style? Turing Award lecture, CACM 21(8), 1978}
  \slideref{J. Backus. \textbf{Can Programming Be Liberated from the von Neumann Style?} CACM 1978. \rfsep
    R. Bird. \textbf{An Introduction to the Theory of Lists}. 1987. \rfsep
    G. Malcolm. \textbf{Data Structures and Program Transformation}. SCP 1990. \rfsep
    E. Meijer, M. Fokkinga, R. Paterson. \textbf{Functional Programming with Bananas, Lenses, Envelopes and Barbed Wire}. FPCA 1991. \rfsep
    R. Bird, O. de Moor. \textbf{Algebra of Programming}. Prentice Hall, 1997. \rfsep
    R. Hinze, N. Wu, J. Gibbons. \textbf{Conjugate Hylomorphisms, or: The Mother of All Structured Recursion Schemes}. POPL 2015. \rfsep
    J. Gibbons. \textbf{Patterns in Functional Programming}. CUP, forthcoming (2027).}
  \speakernote{Backus 1978 to Gibbons 2027 \\ Quote: Backus, Turing lecture}
\end{frame}

\begin{frame}
  \frametitle{Initial Algebras and Catamorphisms}
  \[ \List\,A \;\cong\; 1 + A \times \List\,A \;=\; F\,(\List\,A)
     \qquad \inn = [\op{nil}, \op{cons}] : F\,(\List\,A) \to \List\,A \]
  \[
    \begin{tikzcd}[alt={Commuting square: F of mu F maps to F B by F of
        the fold; in maps F mu F to mu F; a maps F B to B; the fold maps
        mu F to B.}]
      F\,\mu F \arrow[r, "F\,\cata{a}"] \arrow[d, "\inn"'] \& F\,B
        \arrow[d, "a"] \\
      \mu F \arrow[r, "\cata{a}"', dashed] \& B
    \end{tikzcd}
    \qquad
    f = \cata{a} \iff f \circ \inn = a \circ F\,f
  \]
  \begin{itemize}
    \item An $F$-algebra $a : F\,B \to B$; the fold $\cata{a}$ is the
      \emph{unique} algebra morphism out of $(\mu F, \inn)$.
    \item For lists, $\cata{[\op{base}, \op{step}]} = \op{foldr}\;
      \op{step}\; \op{base}$.
  \end{itemize}
  \speakernote{Universal property: every law below is uniqueness \\ F f:
      fmap, recursive positions only}
\end{frame}

\begin{frame}
  \frametitle{Fusion}
  \[
    \begin{tikzcd}[alt={Two commuting squares side by side. Left: the fold
        of a from mu F to B. Right: h from B to C, with F h on top and
        algebras a and b. The outer rectangle shows h after the fold of a
        equals the fold of b.}]
      F\,\mu F \arrow[r, "F\,\cata{a}"] \arrow[d, "\inn"'] \& F\,B
        \arrow[r, "F\,h"] \arrow[d, "a"] \& F\,C \arrow[d, "b"] \\
      \mu F \arrow[r, "\cata{a}"'] \arrow[rr, bend right = 30,
        "\cata{b}"', dashed] \& B \arrow[r, "h"'] \& C
    \end{tikzcd}
  \]
  \[ h \circ a = b \circ F\,h \implies h \circ \cata{a} = \cata{b} \]
  Two traversals become one, if the \textbf{side condition} holds.
  \speakernote{Side condition: h an algebra morphism \\ Proof: outer
      rectangle commutes; uniqueness}
\end{frame}

\begin{frame*}[vertical-alignment = top]
  \frametitle{Hylomorphisms: Divide and Conquer}
  \vspace*{2mm}
  \[
    \begin{tikzcd}[alt={Commuting square: h maps A to B; c maps A to F A;
        F h maps F A to F B; a maps F B to B.}, row sep = 7mm]
      A \arrow[r, "\hylo{a}{c}", dashed] \arrow[d, "c"'] \& B \\
      F\,A \arrow[r, "F\,\hylo{a}{c}"'] \& F\,B \arrow[u, "a"']
    \end{tikzcd}
    \qquad
    \hylo{a}{c} = a \circ F\,\hylo{a}{c} \circ c
  \]
  \hsblock[before upper = {\fvset{fontsize = \footnotesize}}]{hylo-qsort.hs}
  \speakernote{c divides, a conquers \\ Names as in the library: merge,
      split}
\end{frame*}

\begin{frame}[vertical-alignment = top]
  \frametitle{The Call Tree of Quicksort}
  \vspace*{5mm}
  \begin{center}
    \begin{tikzpicture}[alt={The call tree of quicksort on the list 3, 1,
        4, 5, 2, drawn as a tree of layers. The root is Node 3, for the
        input 3, 1, 4, 5, 2. Its children are Node 1, for the input 1, 2,
        and Node 4, for 4, 5. Node 1 has children Leaf and Node 2; Node 4
        has children Leaf and Node 5; Node 2 and Node 5 have two Leaf
        children each. The coalgebra split builds the tree from the root
        down; the algebra merge folds it from the leaves up.},
      level distance = 8.5mm,
      level 1/.style = {sibling distance = 52mm},
      level 2/.style = {sibling distance = 26mm},
      level 3/.style = {sibling distance = 13mm},
      every node/.style = {box, font = \small\ttfamily},
      leaf/.style = {font = \scriptsize\ttfamily, draw = black!30,
        fill = black!5},
      inp/.style = {draw = none, fill = none, font = \scriptsize,
        text = gray, left = 1.5mm}]
      \node (r) {Node 3}
        child {node (l) {Node 1}
          child {node[leaf] {Leaf}}
          child {node (l2) {Node 2}
            child {node[leaf] {Leaf}}
            child {node[leaf] {Leaf}}}}
        child {node (rr) {Node 4}
          child {node[leaf] {Leaf}}
          child {node (r2) {Node 5}
            child {node[leaf] {Leaf}}
            child {node[leaf] {Leaf}}}};
      \node[inp] at (r.west) {$[3,1,4,5,2]$};
      \node[inp] at (l.west) {$[1,2]$};
      \node[inp] at (rr.west) {$[4,5]$};
      \node[inp] at (l2.west) {$[2]$};
      \node[inp] at (r2.west) {$[5]$};
      \draw[->, very thick, accent] (-6.0, 0.2) -- node[left, draw = none,
        fill = none, note, text = accent] {$\ana{\code{split}}$} (-6.0, -2.8);
      \draw[->, very thick, accent2] (6.0, -2.8) -- node[right, draw = none,
        fill = none, note] {$\cata{\code{merge}}$} (6.0, 0.2);
    \end{tikzpicture}
  \end{center}
  \begin{itemize}
    \item A value of $\mu\,\code{TreeF}$: the \textbf{call tree}.
      $\hylo{a}{c} = \cata{a} \circ \ana{c}$, when the tree is finite.
    \item<2-> A hylomorphism consumes each layer as it is produced: the tree
      is \textbf{virtual}.
  \end{itemize}
  \speakernote{Unfold builds it, fold consumes it \\ hylo: never built}
\end{frame}

\begin{frame}
  \frametitle{Recursive Coalgebras and Their Laws}
  $c$ \emph{recursive}: a unique solution $\hylo{a}{c}$ for every $a$.
  \[
    \begin{tikzcd}[alt={Left fusion: f after the hylo of a and c equals the
        hylo of a prime and c.}]
      A \arrow[r, "\hylo{a}{c}"] \arrow[rd, "\hylo{a'}{c}"', dashed] \& B
        \arrow[d, "f"] \\
      \& C
    \end{tikzcd}
    \quad
    \begin{tikzcd}[alt={Deforestation: the unfold of c into mu F followed
        by the fold of b equals the hylo of b and c.}]
      A \arrow[r, "\ana{c}"] \arrow[rr, bend right = 30, "\hylo{b}{c}"',
        dashed] \& \mu F \arrow[r, "\cata{b}"] \& B
    \end{tikzcd}
  \]
  \[
    \begin{array}{rcl@{\qquad}l}
      f \circ a = a' \circ F\,f & \implies & f \circ \hylo{a}{c} =
        \hylo{a'}{c} & \text{(fusion, left)}\\
      c \circ f = F\,f \circ c' & \implies & \hylo{a}{c} \circ f =
        \hylo{a}{c'} & \text{(fusion, right)}\\
      c' \circ a = \op{id} & \implies & \hylo{b}{c'} \circ \hylo{a}{c} =
        \hylo{b}{c} & \text{(deforestation)}
    \end{array}
  \]
  \slideref{V. Capretta, T. Uustalu, V. Vene. \textbf{Recursive Coalgebras from Comonads}. Inf. Comput. 2006. \rfsep
    R. Hinze, N. Wu, J. Gibbons. \textbf{Conjugate Hylomorphisms, or: The Mother of All Structured Recursion Schemes}. POPL 2015.}
  \speakernote{All three from uniqueness \\ Deforestation: no intermediate
      structure}
\end{frame}

\begin{frame*}[vertical-alignment = top]
  \frametitle{Calculating Sort-Then-Double}
  \vspace*{4mm}
  \hsblock{qsort-calc.hs}
  \begin{onlyenv}<2->
    The \textbf{side condition} of this fusion, checked here only
    \textbf{by hand}:
    \hsblock{qsort-side.hs}
  \end{onlyenv}
  \par
  \begin{onlyenv}<3->
    \begin{center}
      Two passes and a temporary list $\Longrightarrow$ \textbf{one pass},
      doubling each pivot.
    \end{center}
  \end{onlyenv}
  \par
  \speakernote{Spec: sort, then double \\ Side condition by hand; in Rocq,
      Part III}
\end{frame*}

\begin{frame*}[vertical-alignment = top]
  \frametitle{List Homomorphisms Are Hylomorphisms}
  \vspace*{3mm}
  \hsblock[before upper = {\fvset{fontsize = \scriptsize}}]{hom.hs}
  \begin{itemize}
    \item<2-> $h\,(xs \mathbin{+\!\!+} ys) = h\,xs \odot h\,ys$:
      \textbf{first homomorphism theorem}, $h = \op{reduce}\,(\odot) \circ
      \op{map}\,f$.
    \item<3-> Any cut, same result (induction on the call tree): the cut
      is a \textbf{scheduling} decision.
  \end{itemize}
  \slideref{R. Bird. \textbf{An Introduction to the Theory of Lists}. Logic of Programming and Calculi of Discrete Design, 1987.}
  \speakernote{The cut is a parameter of divide \\ Halves: parallel;
      heads: a loop}
\end{frame*}

\begin{frame*}[vertical-alignment = top]
  \frametitle{The Third Homomorphism Theorem}
  \vspace*{3mm}
  \textbf{Theorem.} If $h$ is both leftwards and rightwards,
  \[ h\,(x : xs) = x \oplus h\,xs \qquad\text{and}\qquad
     h\,(xs \mathbin{+\!\!+} [x]) = h\,xs \otimes x, \]
  then $h$ is a list homomorphism, with
  \hsblock{third-combiner.hs}
  \begin{itemize}
    \item<2-> A loop from the left remembers only $h\,xs$; one from the
      right only $h\,ys$. So $h\,(xs \mathbin{+\!\!+} ys)$ depends only on
      both.
    \item<3-> Insertion sort, from either side $\Longrightarrow$ mergesort,
      by calculation.
    \item<4-> Combiners derived automatically: weak right inverses (Morita
      et al.); synthesised for loops (Farzan and Nicolet).
  \end{itemize}
  \slideref{J. Gibbons. \textbf{The Third Homomorphism Theorem}. JFP 1996. (Conjectured by Bird, proved by Meertens, 1989.) \rfsep
    K. Morita et al. \textbf{Automatic Inversion Generates Divide-and-Conquer Parallel Programs}. PLDI 2007. \rfsep
    A. Farzan, V. Nicolet. \textbf{Synthesis of Divide and Conquer Parallelism for Loops}. PLDI 2017.}
  \speakernote{Two sequential programs agree: a parallel one exists \\ g:
      a right inverse of h}
\end{frame*}

\begin{frame*}[vertical-alignment = top]
  \frametitle{One Specification, Many Parallel Programs}
  \vspace*{3mm}
  \hsblock[before upper = {\fvset{fontsize = \footnotesize}}]{psort.hs}
  \begin{center}
    \begin{tikzpicture}[alt={A three-stage parallel sort. A list enters a
        split stage, which halves every list longer than n and feeds the
        halves back into its own input. The short lists go to a farm of k
        workers that sort them. The sorted lists go to a merge stage, which
        merges two sorted lists and feeds the result back, until one sorted
        list remains.},
      stage/.style = {box, font = \small, minimum height = 7mm,
        minimum width = 12mm, inner sep = 3pt},
      loop/.style = {->, accent2, thick},
      >={Stealth[length = 1.5mm]}]
      \node (in) {list};
      \node[stage, right = 8mm of in] (sp) {split};
      \foreach \i/\y in {1/5, 2/0, 3/-5}
        \node[stage, right = 16mm of sp, yshift = \y mm, minimum height = 5mm]
          (w\i) {sort};
      \node[stage, right = 16mm of w2] (mg) {merge};
      \node[right = 8mm of mg] (out) {sorted};
      \draw[->] (in) -- (sp);
      \foreach \i in {1, 2, 3} { \draw[->] (sp) -- (w\i);
        \draw[->] (w\i) -- (mg); }
      \draw[->] (mg) -- (out);
      \draw[loop] (sp.north east) .. controls +(0.6, 0.8) and +(-0.6, 0.8)
        .. (sp.north west);
      \draw[loop] (mg.north east) .. controls +(0.6, 0.8) and +(-0.6, 0.8)
        .. (mg.north west);
      \node[note, above = 6mm of sp] {halves fed back};
      \node[note, above = 6mm of mg] {merged pairs fed back};
      \node[note, above = 0mm of w1] {farm of $k$};
    \end{tikzpicture}
  \end{center}
  \begin{onlyenv}<2->
    {\centering $\op{psort} : \List\,\op{Int} \xrightarrow{\;\sigma\;}
      \List\,\op{Int}$ \qquad $\sigma = \op{FB}\,(\op{FUN}\,A) \parallel
      \op{FARM}\,k\,(\op{FUN}\,A) \parallel \op{FB}\,(\op{FUN}\,A)$\par}
    \smallskip
    Types carry the parallel structure; type checking: is there an equal
    program with this structure (normalise with the hylo laws, unify)? A
    \textbf{cost model} picks $\sigma$, $n$ and $k$.
  \end{onlyenv}
  \slideref{D. Castro, K. Hammond, S. Sarkar. \textbf{Farms, Pipes, Streams and Reforestation: Reasoning about Structured Parallel Processes using Types and Hylomorphisms}. ICFP 2016. \rfsep
    D. Castro, K. Hammond, S. Sarkar, Y. Alguwaifli. \textbf{Automatically Deriving Cost Models for Structured Parallel Processes Using Hylomorphisms}. FGCS 2018.}
  \icfptag
  \speakernote{Ours, ICFP 2016 \\ Reforestation exposes stages \\ Cost model
      picks sigma, n, k}
\end{frame*}

\begin{frame}[vertical-alignment = top]
  \frametitle{What Was Missing}
  \vspace*{2mm}
  \begin{itemize}[<+->]
    \item \textbf{Laws as axioms}: the ICFP 2016 type system assumed them
      (sound for a cpo semantics); programs could diverge.
    \item \textbf{Side conditions on trust}: \code{hylo merge halves}
      becomes \code{... || fb (fun merge)} only if \code{merge} is
      associative and commutative; the library took the programmer's word
      for it.
    \item \textbf{Unverified tools}: they implement the schemes or the
      rewrites; calculations by hand or by trusted rewriting.
    \item \textbf{No idiomatic code}: calculated programs stay
      higher-order; we trust the compiler to inline them.
  \end{itemize}
  \uncover<5->{\begin{center}
    \begin{tikzpicture}[alt={The question for the rest of the talk: can a
        proof assistant prove each law once, check every rewrite, side
        conditions included, and still produce efficient code?}]
      \node[hot, text width = 140mm, align = center, inner sep = 1.5mm]
        {Can a proof assistant \textbf{prove} each law once,
        \textbf{check} every rewrite, side conditions included, and still
        produce \textbf{efficient} code?};
    \end{tikzpicture}
  \end{center}}
  \slideref{Tools: HYLO (Hu et al., ICFP'96); MAG (de Moor, Sittampalam, TCS'01); Pointless Haskell, DrHylo (Cunha et al., IFL'05); SkeTo (Matsuzaki, Emoto, HLPP'10); HERMIT (Farmer et al., Haskell'12).}
  \icfptag
  \speakernote{Laws assumed; side conditions on trust \\ Same in other tools
      \\ So: a proof assistant}
\end{frame}
